% Phase 7 text for the course template, trimmed for the 10-page report. Each block is headed by the template section it belongs to.
% Add phase7_refs.bib to references.bib. The longer version, with the gold-size figure, is phase7_sections_full.tex.

% ---------- Section 3.2, Corpus statistics ----------
\paragraph{Phase 7.} Train, dev and test hold 827, 118 and 237 queries against separate pools of 4320, 1023 and 1727 candidates. Gold counts are heavily skewed (dev mean 8.2, median 5, maximum 92), so with seven predictions a perfect ranking reaches only 0.618 dev Micro-F1. Documents are long: the median dev query has 4451 words and the 90th percentile 14684. Dev queries also appear in their own pools, so we remove self-matches, and gold cases were always older than the query in the 352 pairs whose header dates parse, which gives a safe date filter for about half of the documents.

% ---------- Section 3.3, Preprocessing ----------
For Phase 7, query and candidate each get a 512-token excerpt: the 96-token opening plus the sentences sharing the most rare word 4-grams with the other document. We did not scrub the text. Citation traces survive the organisers' masking, and matching them directly raised dev F1 from 0.4603 to 0.4747, but that reads answers from masking gaps and improves no model, so we rejected it and never submitted it. The traces remain in the raw text our features read.

% ---------- Section 4, Phase 7 subsection ----------
\subsection{Phase 7 - Prior Case Retrieval (PCR)}
The system has three stages \cite{nogueira2019passage}. A hashed word 4-gram overlap score orders the pool after removing self-matches and later judgments, reaching 0.4603 dev F1 against 0.101 for unigram BM25, because a citing judgment often quotes the cited one. A gradient-boosted ranker \cite{ke2017lightgbm} then rescores the list from 27 features, including overlap computed only from query sentences near a \texttt{<CITATION>} placeholder, which uses where a citation sits, not which case it was. The ranker is not a Phase 1 model. Finally, Phase 1 backbones jointly read the paired excerpts of the top 50 candidates, and a linear head on the last token (GQA) or the token mean (encoder) gives a score, trained with softmax cross-entropy over one gold case and seven negatives from the lexical top 40. The final score is
\begin{equation}
s(q,c) = z_q\big(\mathrm{logit}\,r(q,c)\big) + 0.2\cdot\frac{1}{3}\sum_{m=1}^{3} z_q\big(n_m(q,c)\big),
\end{equation}
where $z_q$ standardises within a query, $r$ is the ranker and $n_m$ are three rerankers. Four ideas failed: the neural score as a ranker feature, retraining on train and dev, four-epoch training (train top-1 0.95, neural dev F1 near 0.32), and averaging more seeds (test 0.4932).

% ---------- Section 5.3, Ablations ----------
\paragraph{Phase 7.} Table~\ref{tab:p7steps} adds one stage at a time, Table~\ref{tab:p7methods} changes only the adaptation method, and Table~\ref{tab:p7controls} changes one training factor. We call a control difference real only if its interval excludes zero and it clearly exceeds the seed spread (s.d. 0.0097 neural F1 over five encoder seeds).

\begin{table}[H]
\centering\small
\caption{Phase 7 pipeline, with the paired dev gain over the previous row.}
\label{tab:p7steps}
\begin{tabular}{lrrrl}
\toprule
System & Dev F1 & Dev hits & Test F1 & Dev gain (95\% interval) \\
\midrule
Word 4-gram overlap & 0.4603 & 414 & 0.4554 & \\
+ learned ranker & 0.4847 & 436 & 0.4790 & +0.0245 [0.0095, 0.0410] \\
+ citing-context features & 0.4947 & 445 & 0.4891 & +0.0100 [$-$0.0010, 0.0214] \\
+ three rerankers (final) & 0.4992 & 449 & 0.4938 & +0.0044 [0.0010, 0.0094] \\
\bottomrule
\end{tabular}
\end{table}

\begin{table}[H]
\centering\small
\caption{Phase 7 adaptation methods. Neural F1 and AUC as in Table~\ref{tab:p7controls}. Blend: dev, citing-context ranker (0.4947) plus $0.2\,z(\text{neural})$. Test: older ranker (0.4790) plus the same term.}
\label{tab:p7methods}
\begin{tabular}{lrrrrrrrr}
\toprule
Method & Params (M) & Trainable & ms/pair & GB & Neural F1 & AUC & Blend & Test \\
\midrule
Encoder, full & 213.9 & 100\% & 4.5 & 12.0 & 0.3391 & 0.771 & 0.4969 & 0.4849 \\
GQA, full (teacher) & 213.9 & 100\% & 4.2 & 12.0 & 0.3168 & 0.756 & 0.4914 & 0.4838 \\
GQA, LoRMA & 214.6 & 0.31\% & 4.5 & 6.2 & 0.3335 & 0.784 & 0.4914 & 0.4867 \\
GQA, KD student & 123.8 & 100\% & 2.1 & 7.4 & 0.2079 & 0.646 & 0.4958 & 0.4849 \\
GQA, IA$^3$ & 214.0 & 0.03\% & 4.7 & 6.8 & 0.3002 & 0.738 & 0.4936 & -- \\
GQA, partial (top 4) & 213.9 & 21.08\% & 4.2 & 3.7 & 0.3168 & 0.762 & 0.4958 & -- \\
GQA, LoRA & 214.8 & 0.40\% & 4.7 & 6.4 & 0.3135 & 0.760 & 0.4958 & -- \\
\bottomrule
\end{tabular}
\end{table}

\begin{table}[H]
\centering\small
\caption{Phase 7 one-factor controls on dev, one seed each. Neural F1: the model alone ranks the lexical top 50 (chance 0.112). Intervals: paired query bootstrap. Hits: added to the ranker at weight 0.2 (ranker alone 445).}
\label{tab:p7controls}
\begin{tabular}{lcccc}
\toprule
Run & Neural F1 & Difference [95\% interval] & AUC & Hits \\
\midrule
Encoder baseline & 0.3391 & -- & 0.771 & 447 \\
Pairwise loss & 0.3324 & $-0.0067$ [$-0.0173$, 0.0035] & 0.765 & 447 \\
Random negatives & 0.4002 & $+0.0611$ [0.0363, 0.0886] & 0.881 & 442 \\
Last-token pooling & 0.3079 & $-0.0311$ [$-0.0501$, $-0.0135$] & 0.747 & 444 \\
Frozen backbone, head only & 0.2379 & $-0.1012$ [$-0.1282$, $-0.0757$] & 0.676 & 447 \\
Random start, no Phase 1 & 0.1946 & $-0.1445$ [$-0.1784$, $-0.1119$] & 0.629 & 444 \\
\midrule
GQA baseline (last token) & 0.3168 & -- & 0.756 & 442 \\
GQA mean pooling & 0.3357 & $+0.0189$ [$-0.0021$, 0.0406] & 0.779 & 444 \\
\bottomrule
\end{tabular}
\end{table}

% ---------- Section 6, Results ----------
\paragraph{Phase 7.} Our best entry scored 0.4938 Micro-F1@7 on test (precision 0.5039, recall 0.4841). Test tracked dev within about 0.006 at every stage, although the blend was chosen from about 15 dev combinations, so the last dev gain is optimistic. Most of the score comes from exact text reuse: the neural term adds only about eight test hits. Neural models are weak alone (AUC 0.65--0.78 in the top 50, against 0.915 for lexical order) but correlate 0.17 or less with lexical order, so a small additive weight helps.

All seven adaptation methods lie within five dev hits once blended. LoRMA \cite{bihany2025lorma} matches full fine-tuning with 0.31\% of the weights and about half the memory; LoRA \cite{hu2022lora} and IA$^3$ \cite{liu2022ia3} behave similarly. We suspect that 827 training queries cannot use the capacity of full fine-tuning, which fits the overfitting in longer runs, but we did not test this. The encoder's lead over GQA is mostly pooling: with mean pooling on both, the scores are 0.3391 and 0.3357. Distillation \cite{hinton2015distilling} cut parameters by 42\% and doubled scoring speed, but lost a third of standalone quality (0.2079 against 0.3168, difference interval $[-0.1385, -0.0806]$). Its blended score does not show that quality survived, only that a cheaper reranker may suffice for so small a correction.

Phase 1 transfer is real for the standalone model: 0.1946 from a random start (learning rate 1e-4), 0.2379 with a frozen backbone and 0.3391 with full fine-tuning. But the frozen model ties the fine-tuned one at 447 blended hits, so the blended gain cannot be credited to fine-tuning. Mean pooling beats last-token pooling on the encoder, and pairwise and listwise losses do not differ. Random negatives gave the largest standalone gain and the worst blend, 442 hits against 447 (interval $[-0.0110, -0.0011]$, uncorrected). Their scores correlate 0.555 with the ranker, against 0.245 for hard negatives, which suggests they teach topical relatedness the ranker already has.

% ---------- Section 7.1, Failure modes ----------
\paragraph{Phase 7.} Of 973 dev gold cases, 449 were retrieved, 417 cannot be retrieved because their queries have more than seven gold cases, and 203 lie outside the top 100 (these sets overlap). Queries with one or two gold cases cannot exceed about 0.2 precision, and those with 21 or more reach 0.98 precision but 0.19 recall: both are metric ceilings. Recall falls from 0.68 to 0.23 as queries grow, but hits over available slots rise from 0.65 to 0.90, so long queries are not harder. Long and very short candidates are: where queries have at most seven gold cases, recall is 0.84 for candidates of 4--8k words, 0.67 above 8k and 0.50 below 1k, although these bands hold only 40--78 pairs.

The main remaining error is citation without shared text. Missed gold cases have a mean overlap z-score of 0.56, against 3.37 for hits, and false positives sit at 1.52. The neural term swapped only 58 of 826 top-7 entries, gaining 21 gold and losing 17; its excerpts are chosen by the same n-gram overlap, which may limit it to evidence the lexical stage already has.

% ---------- Section 7.2, What the metrics do not measure ----------
Gold labels record citation, not legal relevance, so an uncited but relevant precedent counts as an error in Phase 7.

% ---------- Section 7.3, Limitations and next steps ----------
In Phase 7 one hit is worth about 0.0006 test F1, so blend gains are small. Controls use one seed each, and the random start bounds transfer only for this two-epoch budget. Next we would score chunks of each candidate, so the neural stage sees text the lexical stage did not select, and test mixed hard and random negatives.

% ---------- Section 7.4, Surprising findings ----------
In Phase 7, a position-only signal was worth more than any neural model (17 test hits against about eight), and the model with the best standalone score was the worst blend component.

% ---------- Why techniques do or do not transfer (Section 7.1 or 8) ----------
Classification phases predict a label from one document, which pretrained representations can summarise. PCR depends on a relation between two long documents, often verbatim reuse, which exact n-gram matching captures directly and a 1024-token excerpt only partly preserves.
